Una pequeña empresa artesanal fabrica dos tipos de jabones ecológicos: de \textbf{lavanda}  y de \textbf{romero}  Ambos se elaboran a partir de aceites esenciales y glicerina vegetal, y el objetivo de la empresa es \textbf{maximizar el beneficio semanal}.

Cada kilogramo de jabón requiere determinadas cantidades de \textbf{aceite esencial (ml)} y \textbf{glicerina (ml)}, y genera un beneficio unitario según la siguiente tabla:

\[
\begin{array}{lccc}
\text{Producto} & \text{Aceite esencial (ml/kg)} & \text{Glicerina (ml/kg)} & \text{Beneficio (€ / kg)} \\
\hline
\text{Jabón de lavanda } & 40 & 100 & 5 \\
\text{Jabón de romero }  & 60 & 80  & 4 \\
\hline
\end{array}
\]

La empresa dispone semanalmente de \textbf{8\,000 ml de aceite esencial} y \textbf{12\,000 ml de glicerina}. 

Además, por cuestiones de marketing, se deben fabricar al menos \textbf{60 kg en total} de jabón.

Por último, al como máximo el 50\% de la glicerina consumida debe destinarse al jabón de romero.

El siguiente modelo de programación permite obtener el plan de producción óptimo, donde ($x_1$ y $x_2$ representan los kg de jabón de lavanda y de romero, respectivamente):

\[
\begin{aligned}
\mbox{max.} \quad & z = 5x_1 + 4x_2 \\
\text{s.a.} \quad 
& 40x_1 + 60x_2 \le 8\,000 & & \text{(disponibilidad de aceite esencial)} \\
& 100x_1 + 80x_2 \le 12\,000 & & \text{(disponibilidad de glicerina)} \\
& x_1 + x_2 \ge 60 & & \text{(demanda mínima total)} \\
& 5x_1 - 4x_2 \ge 0 & & \text{(glicerina del romero)} \\
& x_1, x_2 \ge 0
\end{aligned}
\]

Se conoce la tabla correspondiente a la solución óptima del problema:

% \begin{center}
% \begin{tabular}{c|c|cccccc|}
%  & $z$ & $x_1$ & $x_2$ & $h_1$ & $h_2$ & $h_3$ & $h_4$\\ 
%        & $-600$ & $0$ & $0$ & $0$ & $-1/20$ & $0$ & $0$ \\ 
% \hline
% $h_1$ & 1520  & $0$ & $0$ & $1$ & $-27/50$ & $0$ & $-14/5$\\ 
% $h_3$ & 72  & $0$ & $0$ & $0$ & $11/1000$ & $1$ & $1/50$\\ 
% $x_2$ & 60  & $0$ & $1$ & $0$ & $1/200$ & $0$ & $1/10$\\ 
% $x_1$ & 72  & $1$ & $0$ & $0$ & $3/500$ & $0$ & $-2/25$\\ 
% \hline
% \end{tabular}
% \end{center}

\begin{center}
\begin{tabular}{c|c|cccccc|}
 & $z$ & $x_1$ & $x_2$ & $h_1$ & $h_2$ & $h_3$ & $h_4$\\ 
Phase 2 & $-600$ & $0$ & $0$ & $0$ & $-1/20$ & $0$ & $0$ \\ 
\hline
$h_1$ & 1100  & $0$ & $0$ & $1$ & $-23/40$ & $0$ & $-7/2$\\ 
$h_3$ & 75  & $0$ & $0$ & $0$ & $9/800$ & $1$ & $1/40$\\ 
$x_2$ & 75  & $0$ & $1$ & $0$ & $1/160$ & $0$ & $1/8$\\ 
$x_1$ & 60  & $1$ & $0$ & $0$ & $1/200$ & $0$ & $-1/10$\\ 
\hline
\end{tabular}
\end{center}



Se pide:
\begin{enumerate}[label=\alph*)]
    \item Explicar por qué el requisito ``como máximo el 50\% de la glicerina consumida debe destinarse al jabón de romero'' se formula como $5x_1 - 6x_2 = 0$.

    Si cada kilogramo de lavanda y de romero requiere 100 y 80 ml de glicerina, respectivamente, se cumple:
    
    \[
    80x_2 \le 0.5(100x_1 + 80x_2)
    \;\Longleftrightarrow\;
    5x_1 - 4x_2 \ge 0.
    \]
    

    \item Interpretar la solución óptima en términos del beneficio, el plan de producción y el cumplimiento de los requisitos.

        El plan de producción óptimo consiste en:

    \begin{itemize}
        \item Producir 60 kg de aceite de lavanda y 75 kg de aceite de romero. 
        \item Emplear $7\,900$ ml de aceite esencial de los $8\,000$ disponibles.
        \item Agotar los $12\,000$ ml de glicerina.
        \item Dedicar estrictamente el 50\% de la glicerina a jabón de romero.
    \end{itemize}

    Esto reporta un beneficio semanal de 600€.
    
    \item Identificar la base y la inversa de la base de la solución óptima.

    La base se obtiene a partir de las columnas en la matriz $A$ de las cuatro variables básicas ($h_1$, $h_3$, $x_2$, $x_1$): 

    \begin{equation}
    \begin{split}
    B =\begin{pmatrix}
    1 & 0 & 60 & 40\\
    0 & 0 & 80 & 100\\
    0 & -1 & 1 & 1\\
    0 & 0 & -4 & 5\\
    \end{pmatrix}
    \end{split}
    \end{equation}

    La inversa de la base se obtiene a partir de la tabla correpsondiente a la solución, tomando las columnas donde originalmente estaban las columnas de la matriz identidad o las opuestas de dichas columnas (cambiando su signo):

    \begin{equation}
    \begin{split}
    B^{-1} =\begin{pmatrix}
    1 & -23/40 & 0 & 7/2\\
    0 & 9/800 & -1 & -1/40\\
    0 & 1/160 & 0 & -1/8\\
    0 & 1/200 & 0 & 1/10\\
    \end{pmatrix}
    \end{split}
    \end{equation}

    \item ¿Cuál sería el precio que debería estar dispuesta a aceptar la empresa para renunciar a 1 litro de glicerina en una semana dada?¿Cuántos litros estaría dispuesta a pagar por ese precio?

     El precio sombra la restricción 2 (disponibilidad de glicerina) es $\pi^B_2 = 1/20$, con lo que por cada ml adicional de de glicerina de la que disponga el beneficio aumentará en $0.05$ €.  Como el problema es lineal, si se dispusiera de un litro adicional, el beneficio podría ser de 50€. 
     
     Sin embargo, hay que comprobar el rango de $b_2$ dentro del cual la base de la tabla sigue siendo una base óptima.
     
    \begin{equation}
    \begin{split}
    u^B=B^{-1}b =\begin{pmatrix}
    1 & -23/40 & 0 & 7/2\\
    0 & 9/800 & -1 & -1/40\\
    0 & 1/160 & 0 & -1/8\\
    0 & 1/200 & 0 & 1/10\\
    \end{pmatrix}
    \begin{pmatrix}
        8000\\
        b_2\\
        60\\
        0
    \end{pmatrix}   
    \geq 0 \Rightarrow \cdots \Rightarrow 6000 \leq b_2 \leq 20000
    \end{split}
    \end{equation}

    Es decir, la empresa estaría dispuesta a comprar hasta 8 litros de glicerina por un coste máximo de 50€ por litro.

\end{enumerate}